
\documentclass[a4paper, 12pt]{article}
\usepackage[utf8]{inputenc} \usepackage[T1]{fontenc}
\usepackage{textcomp}
\usepackage{amssymb}
\usepackage{baskervald} \usepackage{newtxmath}
\usepackage{amsmath, amssymb}
\newtheorem{problem}{Problem}
\newtheorem{example}{Example}
\newtheorem{lemma}{Lemma}
\newtheorem{theorem}{Theorem}
\newtheorem{problem}{Problem}
\newtheorem{example}{Example} \newtheorem{definition}{Definition}
\newtheorem{lemma}{Lemma}
\newtheorem{theorem}{Theorem}
\DeclareMathAlphabet{\mathcal}{OMS}{cmsy}{m}{n}
\usepackage{parskip}

\begin{document}

\section*{La línea de sombra}

\begin{itemize}
    \item Todavía no vinculamos nada a Coleridge.
\end{itemize}

Hay tres niveles narrativos: historia de marinos, historia autobiográfica, y
viaje metafórico. 

De acuerdo con Conrad, 

\begin{quote}
    La primera intención de esta obra era el presentar ciertos hechos referentes
    a ese \textit{instante} en que la juventud despreocupada y fogosa alcanza la
    época más consciente y conmovedora de la madurez.
\end{quote}

En ese instante, en ese momento de la vida, uno habita un punto medio cruel: las
ilusiones adolescentes y la ingenuidad, ya desmoronadas, no nos dan fuerza, pero
la placidez y la sabiduría que nos dan los años tampoco han llegado todavía.
Recordar versos de Baudelaire en \textit{Spleen de París}:

\begin{quote}
    \textit{Je suis comme le roi d'un pays pluvieux},\newline
    \textit{Riche, mais impuissant, jeune et pourtant très-vieux}.
\end{quote}

La frontera que separa la juventud de la adultez es precisamente metaforizada
como una línea de sombra, que describe la vaga experiencia subjetiva, pero a la
vez universal, de madurar a travez de la experiencia y la adversidad. 


\begin{quote}
    Uno cierra tras de sí la pequeña verja de la infancia
    [\textit{boyishness}]---y entra en un jardín encantado. Hasta sus sombras
    resplandecen de promesa. Cada vuelta del camino tiene su seducción. Y no es
    porque se trate de un país inexploraod. Uno sabe muy bien que toda la
    humanidad ha desfilado ese camino. Es el encanto de la experiencia
    universal, en la cual uno espera una sensación excepcional o personal---algo
    de uno mismo.

    Uno va reconociendo las marcas dejadas por sus predecesores, entusiasmado,
    entretenido, tomando la buena y la mala suerte por igual---las patadas y los
    medio centavos, como dice el dicho---la pintoresca suerte común que guarda
    tantas posibilidades para quien las merece o quizás para los afortunados.
    Sí. Uno va. Y el tiempo, también, se va---hasta que uno percibe más allá una
    línea de sombra que le advierte que esa región de juventud temprana,
    también, debe dejarse atrás.
\end{quote}

El propio Conrad la llamó «\textit{a fairly complex piece of work}».
\textit{Bullet point: maturation through hardship}. \textit{From his biography:}
The \textit{Shadow Line} is almost entirely atmosphere, and is poetry as
Coleridge's \textit{Ancient Mariner} was poetry.

Desde una perspectiva histórica, el tema no es tan abstracto como parece: fue
escrito en 1916, durante la primera guerra mundial, en la que el propio hijo de
Conrad participó. Conrad, que llamó a esta guerra la «prueba de una generación»,
ofrece al mundo literario un examen autobiográfico de su propia transición a la
adultez, por distante que sea la experienca de un marinero de la de un soldado.
El elemento autobiográfico no es la nostalgia de un hombre retirado, sino una
respuesta al eterno problema, ahora enfrentado por una nueva generación, de cómo
encontrar sentido o alcanzar un ideal en el arduo paso de la ingenuidad a la
vida. 

La historia comienza con un motivio inicial típico: la \textit{pérdida de
status}. El protagonista abandona, casi caprichosamente, una posición segura en
un contexto que le es familiar, sintiendo que no hay nuevas verdades que
descubrir allí. Esta renuncia no es un acto de arrojo valiente a lo desconocido,
sino una reacción impulsiva ante la incertitud insatisfecha que lo consume.
Desembarca en un puerto oriental y se hospeda en un \textit{Hogar de Oficiales
Marineros}, un centro de alojamiento típico en ciuddes portuarias. Sentimos la
inquietud y la soledad del protagonista, que acaba de abandonarlo todo, y se
predispone a regresar a Inglaterra en un barco que arribará en pocos días.
Contra esta inquietud hay un mundo burocrático, frío e indiferente. 

En el hogar de oficiales marineros se encontrará con dos personajes
determinantes: el capitán Giles, experimentado, sencillo y noble, y un
insoportable llamado Hamilton. Ambos van a representar algo significativo para
nuestro personaje: Hamilton será 

\begin{quote}
the first instance of harm being attempted to be done to me—at any rate,
the first I had ever found out. And I was still young enough, still too much on
this side of the shadow line, not to be surprised and indignant at such things.
\end{quote}

El capitán Giles, por su parte, comprendiendo el punto de inflexión en que se
encuentra el protagonista, facilita ---no diremos cómo--- que el protagonista se
embarque en un viaje. Lo hace prometiéndole el \textit{comando} de un nuevo
barco, y esta palabra, \textit{comando}, es la llave mágica para convencerlo:

\begin{quote}
    All at once, as if a page of a book had been turned over disclosing a word
    which made plain all that had gone before, I perceived that this matter had
    also another than an ethical aspect.

And still I did not move. Captain Giles lost his patience a little. With an
angry puff at his pipe he turned his back on my hesitation.

But it was not hesitation on my part. I had been, if I may express myself so,
put out of gear mentally. But as soon as I had convinced myself that this stale,
unprofitable world of my discontent contained such a thing as a command to be
seized, I recovered my powers of locomotion. 
\end{quote}

El barco que debe comandar está varado en Bankgok: su capitán anterior murió y
fue enterrado en altamar. El protagonista espear como un caballero que abandona
el mundo conocido para rescatar a una doncella. Su meta espiritual
---semi-consciente--- es el autoconocimiento y penetrar en el sentido de la vida
y de un mundo donde Dios es inexistente o está oculto y no existen verdades
trascendentales. El capitán anterior es la cara negativa del protagonista: por
un lado era un misántropo cruel que condenó a su propia tripulación a la ruina;
por otro, es el antecesor del protagonista en una «dinastía» de capitanes que no
se hereda por sangre sino por experiencia y sentido del deber. 



Como género ocupa un lugar intermedio entre la autobiografía y la novela de
aprendizaje. Y en términos de narrativa es doblemente romántica y realista.
Romántica en el sentido de \textit{romance}: una narrativa en este caso secular,
en el sentido de que el héroe no es una figura mítica o divina, caracterizada
por movimientos de descenso y ascenso. El descenso lleva al héroe a niveles cada
vez más bajos de existencia o de consciencia, terminando con su muerte
(metafórica o literal). El ascenso es la anástasis o el retorno. 


De hecho, el carácter romántico de la historia se hace aparente apenas el
personaje asume el comando y empieza, por lo tanto, a considerar la ética
(todavía abstracta) de su deber. Siente repentinamente que habita un sueño: su
nuevo rol contiene un elemento de misticismo que, en cierto sentido, empieza ya,
aunque de modo apenas latente, a revivirlo. Si esta romantización se corresponde
todavía con su inexperiencia o ingenueidad, no lo sabemos.

La doble cara de la narrativa, la tensión entre la línea romántica y la
realista, se refleja precisamente en el lenguaje utilizado por Conrad, que se
mueve constantemente desde la tangibilidad de la vida náutica al misticismo de
la experiencia humana. A frases y palabras como «cuadrante, orientar las vergas,
brioles, escotilla, masteleros», se le oponen adjetivaciones preciosas como:

\begin{quote}
    hechizado, horror fantástico, espectro malicioso, horripilante, hechizo
    maligno, maldición, fantasmas, esqueleto reanimado, terror inconcebible,
    misterio inexpresable, aspecto siniestro, aspecto luminoso, figuras
    sombrías, profundidades inconcebibles, diablería misteriosa, sustancia
    sobrenatural, susurro misterioso
\end{quote}

Hay que comprender que, en 1916, Conrad probablemente sentía que el mundo en el
que se había desarrollado como marinero y escritor estaba agonizando, que la
cultura victoriana que él representó se desvanecía como otra de las muertes
exigidas por la Gran Guerra. En una historia que explora la necesidad de
auto-afirmación, parece traslucirse su necesidad de unir dos puntos de su vida:
el hombre de acción que fue protagonista de la vida, y el hombre de letras que
medita su sentido. \textit{Theme}: the need for answers and directions to deal
with a moment of crisis.


Este es el encanto de las obras de Conrad. Porque, ¿qué las hace más que meras
historias de marinos? Son «nightmarish parables that cannot be deciphered», al
menos no por completo. El propio Conrad dijo en una carta a Sir Sidney Colvin:

\begin{quote}
     Perhaps you won’t find it presumptuous if after 22 years of work I may say
     that I have not been very well understood. I have been called a writer of
     the sea, of the tropics, a descriptive writer, a romantic writer—and also a
     realist. But as a matter of fact all my concern has been with the “ideal”
     values of things, events and people. That and nothing else.

     [Tal vez no le resulte presuntuoso si, después de 22 años de trabajo, me
     permito decir que no he sido muy bien comprendido. Se me ha calificado como
     un escritor del mar, de los trópicos, un escritor descriptivo, un escritor
     romántico... y también como un realista. Pero, en realidad, toda mi
     preocupación ha girado en torno a los valores «ideales» de las cosas, los
     acontecimientos y las personas. Eso y nada más.]
\end{quote}

Estas parábolas están atadas a la vida de Conrad: casi todas sus obras son
experiencias personales dramatizadas con el fin de dotarlas de un sentido
coherente. Pero \textit{Shadow Line}, habiendo sido escrita después de las
novelas más celebres de Conrad (\textit{Lord Jim}, \textit{Heart of Darkness}),
contiene toda la experiencia presente en las obras anteriores. Es la
\textit{auto-examinación final} en una larga cadena de dramatizaciones que
Conrad hizo de su propia vida.

Precisamente, el logro de Conrad como escritor es que ordenó el caos de su
existencia en un arte estructurado que reflejaba su existencia y las realidades
que enfrentaba. Su experiencia, como hombre y escritor, es única en la
literatura inglesa: ninguna expatriación fue tan completa o compleja como la
suya, y tal vez ninguna producción literaria tan profundamente extraña y
creativa.

Un punto notable: es evidente que para Conrad la maduración humana no consiste
en la adquisición de éxito económico ni poder, sino en la conquista de uno
mismo, en el auto-conocimiento y la auto-consciencia. De hecho, como señaló
Edward said, \textit{La línea de sombra} comienza con una admisión de derrota o
fracaso precisamente en el momento de máxima juventud y vitalidad. A partir de
esa derrota, la trama sigue un camino teleológico desde la crisis hasta el
dominio de un mismo. Es una historia de iniciación, de un rito de paso, donde
hacerse adulto es presentado como un acto inconformista de autoafirmación antes
que como una aceptación dócil de las reglaes y roles sociales. 



Pero insistamos: este proceso, si bien tiznado de autobiografía, \textit{no} es
individual. La trama trasciende lo biográfico ampliamente y conforma una ficción
que idealiza la experiencia individual como un patrón universal significativo.
Algo que, por cierto, es muy propio del modernismo literario, donde la reflexión
autobiográfica no versa tanto sobre la vida entera sino sobre periodos
relativamente cortos pero que se consideran cargados de sentido y poder
transformativo.

\textit{Vinculado al párrafo anterior}: \textit{art as a means to confer meaning
to one owns life}. More no this, but what to say?

Estilo más conciso y directo que \textit{Heart of Darkness}. El final de la
novel es inambigüo, la trama es simple y directa, traicionando al modernismo
típico de sus obras anteriores y enfatizando el elemento «tradicionalismo» de la
historia, en el que la capitanía de un barco es vista como una «dynasty,
continuous not in blood, indeed, but in its experience, in its ?, in its
conception of duty».

¿Vale la pena hablar de \textit{delayed decoding}? Otra vez academicismo
desagradable. Delayed decoding $:=$ recurso literario en que un evento se
presenta primero tal y como es percibido por los personajes y sólo luego como es
interpretado, donde existe una diferencia entre la percepción inicial y la
interpretación final (correcta). Supongo que la mención de delayed decoding es
útil sólo para remarcar una diferencia con \textit{Heart of Darkness}: en ésta,
eventos que son inicialmente interpretados como positivos se descubren como
negativos, mientras el proceso en \textit{Shadow Line} es inverso. En
\textit{Heart of Darkness} lo natural adquiere rasgos sobrenaturales; en
\textit{Shadow Line}, a medida que la madurez conquista la adolescencia, las
causas sobrenaturales son sustituidas por lo natural.

\textit{Contaminación, racismo}. .Los británicos también temían "descivilizarse"
debido a su contacto imperial con territorios remotos. El académico Robert
McGill sostiene que la literatura de Joseph Conrad retrata el tenso y
perturbador encuentro del colonialismo con tierras extranjeras, lugares que la
cultura victoriana asociaba profundamente con la enfermedad y la contaminación.
Por ejemplo, en su novela La línea de sombra (ambientada también en el
archipiélago malayo), el narrador encuentra una fotografía de un capitán de
barco junto a una mujer con un atuendo "semi-oriental" y relata, entre
paréntesis, que él "(...incluso la tiró por la borda más tarde)".

Según McGill, estos signos de puntuación funcionan como una suerte de
"cuarentena textual"; una barrera que restringe y aísla el contacto del lector
con "cualquier tipo de representación de lo 'incivilizado'".

Para entender completamente la profundidad de este fragmento (basado en los
análisis de Robert McGill sobre la obra de Conrad), es útil desglosar los tres
conceptos centrales que operan en esta escena:

El miedo a la "Descivilización" (Degeneración moral): Durante el Imperio
Británico, existía una fuerte ansiedad social basada en la idea de que la
identidad "civilizada" europea era frágil. No solo temían a las enfermedades
tropicales reales (como la malaria o el cólera), sino que creían que el contacto
prolongado con culturas no europeas provocaría un "contagio moral". La
fotografía de un capitán británico con una mujer "semi-oriental" representa
exactamente este tabú: la mezcla cultural y racial que los victorianos percibían
como una amenaza a su pureza y civilidad.

El acto del narrador de arrojar la fotografía al mar es un intento desesperado
de purificar el espacio. El barco representa un microcosmos de la sociedad
británica: ordenado, jerárquico y "limpio". Al deshacerse de la foto, el
narrador intenta eliminar la evidencia de la transgresión del capitán anterior y
proteger al barco de esa "contaminación" extranjera.



\pagebreak 
From: https://www.jstor.org/stable/30053144?mag=joseph-conrads-travel-stories-werent-black-and-white&seq=15

Conrad's fictional characters embody the opinions and prejudices to be found in
historical records of colonial life along the coasts of Borneo. The ongoing
rivalry between the colonial powers is reflected in archival documents and
newspaper accounts in which the Dutch regularly blame the English for supporting the illegal
traffic in arms and gunpowder (and slaves), while the English blame the Dutch
for their failure to enforce their own treaty commitments with respect to
slavery. In 1862, William Lingard's trade rival, P. Van Hartrop, complained to
the Dutch authorities that Lingard's provision of rice to famine-stricken
settlers was driving many of them into debt bondage; according to Van Hartrop,
Lingard set ``a price that the small man could not pay, and which led many of
them to fall into slavery with their wives and children'' (qtd. in Campo 92).
In 1876, after being attacked by pirates, the real Lingard published an angry
letter in the \textit{Makassaarc Handelsblad} complaining that Dutch efforts to
guard the coast were ineffective (Warren, ``Fiction'' 26--27; a Campo 105).
This mutual distrust is reflected in the Malay trilogy when Mr. Travers, in
\textit{The Rescue}, wants to ``expose'' the Dutch colonial system (34), or
when the visiting Dutch officers accuse Almayer (rightly) of trading in
contraband gunpowder (\textit{Almayer} 93). In a letter dated September 1897 to
his publisher William Blackwood, Conrad indicated that he was familiar with
gun-running plots like those described in \textit{The Rescue} and
\textit{Almayer's Folly}:

\begin{quote}
In 1848 an Englishman called Wyndham had been living for many years with the
Sultan of Sulu and was the general purveyor of arms and gunpowder. In 1850 or 51
he financed a very lively row in Celebes. He is mentioned in Dutch official
documents as a great nuisance---which he, no doubt, was. I've heard several
versions of his end (occurred in the sixties) all very lamentable. In the 70ies
Lingard had a great if occult influence with the Rajah of Bali. He was a meddler
but very disinterested and was greatly respected by the natives. As late as 1888
arms have been landed on the coast of that island---that to my personal
knowledge. (\textit{Collected Letters} 1: 382)
\end{quote}

In \textit{Conrad's Eastern World} (1966), Norman Sherry cites an apparently
anonymous report published in the Singapore weekly \textit{Straits Times
Overland Journal} for 26 March 1883 (more than four years before Conrad visited
the coast) describing conditions in what its nameless author calls ``the state
of Berouw,'' or in other words Conrad's Sambir:

\begin{quote}
Labour is for women and slaves only. Slaves are met with in almost every house.
On the lower river, there is even a large village wholly inhabited by slaves.
The authorities allow this, in spite of Art. 115 of the Government reg. whereby
slavery in Netherlands India has been abolished. Most of the slaves are fairly
well off excepting those who have to work in the [coal] mines. The number of
these unfortunates yearly sold at Gunong Thabor is estimated at 300. These
people are bought in or kidnapped from the islands of Sooloo and the other
Philippines, and then bartered for gunpowder, muskets, revolvers, lillas
[?], cloths, calico, opium, Dutch candles etc. (130)
\end{quote}

In 1894, seven years after Conrad had left the coast, the slave trade continued,
as a Dutch report confirmed:

























\end{document}



